\section{What the Church Confesses: The Beatific Vision}

The Catholic Creed does not close on a vague wish for immortality. In the Apostles' Creed the believer confesses \textit{carnis resurrectionem, vitam aeternam}---``the resurrection of the body, and life everlasting.'' Every Sunday the Nicene Creed seals the profession of faith with expectation: \textit{Et exspecto resurrectionem mortuorum, et vitam venturi saeculi}---``And I look for the resurrection of the dead, and the life of the world to come.''

Instructing his Lenten catechumens near the Holy Sepulchre, St.\ Cyril of Jerusalem explains why ``life everlasting'' crowns the baptismal symbol: ``The real and true life then is the Father, who through the Son in the Holy Spirit pours forth as from a fountain His heavenly gifts to all; and through His love to man, the blessings of the life eternal are promised without fail to us men also.''\endnote{St.\ Cyril of Jerusalem, \textit{Catechetical Lectures} 18.29, trans.\ Edwin Hamilton Gifford, NPNF2 vol.\ 7.} Eternal life is not bare endless duration---the damned endure forever---but communion in the life of the Triune God.

When heresy or controversy threatened that confession, ecumenical councils and Roman Pontiffs defined the faith with unmistakable precision. In 1215, confronting Cathar dualism, the Fourth Lateran Council promulgated \textit{Firmiter credimus}, confessing that at the end of time Christ will raise all men with their own bodies which they now bear, that they may receive according to their works, whether good or evil---the wicked perpetual punishment with the devil, and the elect everlasting glory with Christ (\textit{qui omnes cum suis propriis resurgent corporibus, quae nunc gestant, ut recipiant secundum opera sua, sive bona fuerint sive mala, illi cum diabolo poenam perpetuam, et isti cum Christo gloriam sempiternam}).\endnote{Fourth Lateran Council (1215), Constitution 1 (\textit{Firmiter credimus}), Denzinger--Sch\"{o}nmetzer (DS) 801.} Near the close of the constitution, the council guards the breadth of heavenly calling: not only virgins and the continent, but also married persons, pleasing God by right faith and good works, merit to arrive at eternal beatitude (\textit{Non solum autem virgines et continentes, verum etiam coniugati, per rectam fidem et operationem bonam placentes Deo, ad aeternam merentur beatitudinem pervenire}).

A century later, a grave dispute arose inside the Latin Church herself. Between 1331 and 1334, Pope John XXII preached at Avignon---expressly as a private theologian, not exercising supreme defining authority---that the souls of the just do not behold the divine essence face to face until after the general resurrection and Last Judgment, contemplating meanwhile only Christ's glorified humanity. Masters at Paris, Franciscan and Dominican theologians, and bishops across Europe appealed against this opinion to Scripture, the Fathers, and the Second Council of Lyons (1274). On his deathbed in December 1334, John XXII retracted his private view and submitted it to the Church's judgment.

Thirteen months later, his successor, Pope Benedict XII, settled the matter forever by a solemn \textit{ex cathedra} definition in the Apostolic Constitution \textit{Benedictus Deus} (29 January 1336).\endnote{Pope Benedict XII, Apostolic Constitution \textit{Benedictus Deus} (29 January 1336), DS 1000--1001.} Exercising apostolic authority, he defined that the souls of the righteous who died before Christ's Passion, together with all the faithful who depart in grace needing no purification---or once fully purified after death---and baptized infants dying before the use of reason, enter immediately after death and any required purification into heaven, the kingdom of heaven, and celestial paradise with Christ, joined to the fellowship of the holy angels, even before the resumption of their bodies and the General Judgment (DS 1000).

Benedict XII defined with equal rigor the very substance of that heavenly life. Since the Passion and death of our Lord Jesus Christ, these blessed souls directly behold the Divine Essence itself in an intuitive and face-to-face vision (\textit{visio intuitiva et facialis}), with no creature standing between them and God as the object beheld; rather, the Divine Essence manifests itself to them immediately, unveiled, brightly, and openly. In thus seeing God, they attain joyful possession and restful enjoyment (\textit{fruitio}) of that very Divine Essence, and by reason of this vision and enjoyment the souls of the blessed already possess true beatitude, eternal life, and rest (DS 1000).

Every precision in the definition shuts out evasion: heavenly knowledge is direct sight rather than deduction from created effects, immediate rather than mediated through any created likeness, and open rather than veiled, issuing directly in \textit{fruitio}, the restful enjoyment of God. In the immediately following paragraph (DS 1001), Benedict XII defined two further consequences of this vision. First, it brings to an end the acts of faith and hope in the blessed, precisely insofar as faith and hope are theological virtues: because faith assents to what is unseen (Heb 11:1) and hope reaches toward what is not yet held (Rom 8:24), clear vision replaces faith and possession fulfills hope, while charity abides forever (1 Cor 13:13). Second, once begun, this vision and enjoyment endure uninterruptedly, without cessation or departure, through the Last Judgment and into eternity.

A century later, in the Bull of Union with the Greeks, \textit{Laetentur caeli} (6 July 1439), the Council of Florence solemnly confirmed Benedict XII and added two further precisions: completely purified souls are received straightway into heaven and clearly behold the Triune and One God Himself as He is, yet with unequal perfection proportioned to the diversity of each saint's merits.\endnote{Council of Florence, Bull of Union with the Greeks \textit{Laetentur caeli} (6 July 1439), DS 1305.} The Object seen is expressly confessed as the Holy Trinity, and the clarity of vision varies according to merit.

After Trent, the \textit{Roman Catechism} gathered this teaching into its exposition of Creed Article XII, ``Life Everlasting.''\endnote{\textit{Catechism of the Council of Trent}, Part I, Creed Article XII, trans.\ J.\ Donovan (Baltimore: Lucas Brothers, 1829), pp.\ 92--97.} The Apostles chose the title \textit{life everlasting} rather than bare \textit{blessedness}, the Catechism explains, to rule out perishable earthly pleasures, to signify that heavenly bliss can never be lost---``Happiness is an accumulation of good without admixture of evil, which, as it fills up the measure of man's desires, must be eternal''---and to acknowledge that no human word equals the reality. Even Scripture's rich titles---Kingdom of God, Paradise, Holy City, New Jerusalem, Father's house---fall short of it.

Citing the Fathers' classical summary, ``an exemption from all evil, and an enjoyment of all good'' (cf.\ Apoc 7:16; 21:4), the \textit{Roman Catechism} distinguishes \textit{essential} from \textit{accessory} beatitude: ``Solid happiness, which we may designate by the common appellation, `essential,' consists in the vision of God, and the enjoyment of his eternal beauty who is the source and principle of all goodness and perfection.'' Accessory happiness embraces the companion gifts surrounding that vision: the glory of the risen body and the radiant fellowship of the saints.

The \textit{Catechism of the Catholic Church} sets this same confession before the contemporary faithful (CCC 1023--1029).\endnote{\textit{Catechism of the Catholic Church}, 2nd ed.\ (1997), \S\S 1023--1029.} Summarizing \textit{Benedictus Deus} and \textit{Laetentur caeli}, it teaches that those who die in God's grace and friendship and are perfectly purified live forever with Christ, beholding God as He is, face to face (CCC 1023). Heaven is therefore no solitary abstraction, but perfect communion of life and love with the Most Holy Trinity, the Blessed Virgin Mary, the angels, and all the blessed---the ultimate end and supreme, definitive fulfillment of every deepest human longing (CCC 1024).

Far from dissolving human personality into an impersonal abyss, heavenly life means dwelling with Christ and in Christ, where the elect do not lose their individuality but discover and retain their true identity and their own name in Him (CCC 1025). Opened to mankind by Christ's death and Resurrection (CCC 1026), this mystery surpasses every human concept and description, so Scripture gestures toward it through converging images---life, light, peace, wedding feast, Father's house, heavenly Jerusalem, and paradise (CCC 1027). Because God transcends every created mind, no creature can behold Him unless God Himself reveals His mystery and grants the supernatural capacity to see Him (CCC 1028). Yet eternal rest is never inert: in heavenly glory the blessed joyfully continue to carry out God's will in relation to other men and to all creation, reigning forever with Christ (CCC 1029; cf.\ Apoc 22:5).
